\documentclass[UTF8,a4paper,12pt]{ctexart}
\usepackage{geometry,amsmath,amssymb,booktabs,longtable,graphicx,float,caption,fancyhdr,setspace,array,enumitem,hyperref,xcolor}
\geometry{left=2.5cm,right=2.5cm,top=2.6cm,bottom=2.5cm}
\setstretch{1.42}
\setlength{\headheight}{15pt}
\pagestyle{fancy}\fancyhf{}\lhead{2015A 太阳影子定位}\rhead{数学建模论文}\cfoot{\thepage}
\hypersetup{colorlinks=true,linkcolor=black,urlcolor=blue}
\captionsetup[table]{labelsep=space}
\captionsetup[figure]{labelsep=space}
\newcommand{\ud}{\mathrm{d}}
\newcommand{\figpath}{../图表/}
\title{基于太阳位置矢量与非线性最小二乘反演的太阳影子定位模型}
\author{}
\date{}
\begin{document}
\maketitle
\begin{abstract}
太阳影子定位问题要求利用直杆影子随时间的变化规律，反推出拍摄地点乃至拍摄日期。本文从太阳赤纬、时角和地平坐标系出发，建立了“时间--太阳高度角/方位角--影尖二维坐标”的机理模型，并在未知坐标轴方向和未知杆高条件下，引入平面旋转角与杆高作为待估参数，用非线性最小二乘完成地点、日期和轨迹的联合反演。

针对问题一，本文建立太阳高度角模型 $L=H/\tan\alpha$，计算天安门广场（北纬 $39^\circ54'26''$、东经 $116^\circ23'29''$）在 2015 年 10 月 22 日北京时间 9:00--15:00 的三米直杆影长。结果表明，影长先减小后增大，9:00 影长为 6.8841 m，15:00 影长为 7.0365 m，最短影长出现在 11:58，约为 3.8411 m。

针对问题二，附件1给出了 2015 年 4 月 18 日某直杆影尖坐标。本文以经纬度、杆高和坐标系旋转角为未知量拟合完整二维轨迹，得到一个主要可能地点为北纬 18.3361$^\circ$、东经 109.7018$^\circ$，拟合杆高 1.9640 m，坐标旋转角 14.9293$^\circ$，二维坐标均方根残差仅 0.000251 m，最大残差 0.000790 m。Bootstrap 扰动下纬度标准差约 0.0359$^\circ$、经度标准差约 0.0945$^\circ$，说明附件1定位结果在观测误差尺度内较稳定。

针对问题三，附件2和附件3日期未知。本文采用全年日期网格搜索与多初值最小二乘联合反演。附件2的最优候选为 2015 年 5 月 23 日、北纬 39.7296$^\circ$、东经 78.8739$^\circ$，RMSE 为 $3.4\times10^{-5}$ m；由于太阳赤纬具有近似对称性，2015 年 7 月 20 日、北纬 39.8040$^\circ$、东经 81.2112$^\circ$也是等价优良候选。附件3的最优候选为 2015 年 11 月 16 日、北纬 30.4121$^\circ$、东经 106.5854$^\circ$，RMSE 为 0.000204 m；2015 年 1 月 24 日、北纬 30.2418$^\circ$、东经 113.2011$^\circ$也是近似等价候选。

针对问题四，当前题目目录中仅有附件4下载说明而无视频本体，自动访问赛题站失败且百度网盘下载需要交互验证。本文不伪造视频测量结果，而给出从视频抽帧、杆底和影尖检测、像素--米标定到地点/日期反演的完整可执行模型。若补充视频文件，可直接用该流程生成影尖坐标并代入本文的联合反演框架。模型检验表明，本文模型具有明确物理含义、残差可追踪、候选解可解释的优点；主要局限在于太阳位置公式近似、影尖识别误差和未知日期下多解性。

\textbf{关键词：}太阳影子定位；太阳高度角；非线性最小二乘；日期网格搜索；视频测量
\end{abstract}
\newpage
\tableofcontents
\newpage

\section{问题重述}
\subsection{问题背景}
视频数据中的太阳影子包含了拍摄时刻、太阳方位和观测地点之间的几何信息。若已知一根直杆竖直立在水平地面上，则杆顶、杆底和影尖构成由太阳光线决定的投影关系。太阳在天空中的位置又由观测地点经纬度、日期和地方时共同决定。因此，影子长度和影尖轨迹能够反向约束地点和日期。这类问题在视频取证、地理定位和图像时空校验中都有实际意义。

本题围绕一根固定直杆的影子变化提出四个任务。第一问要求建立影长随参数变化的数学模型，并计算指定地点、日期和杆高下的影长曲线。第二问给定日期和影尖坐标，要求确定直杆所在地。第三问日期未知，需要同时确定地点和日期。第四问要求利用视频和已知杆高确定拍摄地点，并进一步讨论日期未知时能否确定地点与日期。

\subsection{问题要求与输出目标}
为避免模型输出与题目要求脱节，本文先列出各问的“题目要求--模型输出”映射。
\begin{enumerate}[label=\textbf{问题\arabic*：}]
  \item 题目要求建立影子长度变化模型并绘制天安门三米直杆影长曲线；本文输出太阳高度角公式、北京时间 9:00--15:00 的逐分钟影长表、曲线图和极值时间。
  \item 题目要求由附件1的影尖坐标确定地点；本文输出经纬度候选、拟合杆高、坐标旋转角、残差表和稳定性分析。
  \item 题目要求由附件2、附件3的影尖坐标同时确定地点和日期；本文输出每个附件的日期--地点候选组合、RMSE 排序、轨迹拟合图和多解解释。
  \item 题目要求由视频确定拍摄地点并讨论日期未知情况；本文输出视频处理到影尖坐标的完整建模流程，以及视频缺失条件下的可执行反演方案说明。
\end{enumerate}

\section{问题分析}
\subsection{总体思路}
太阳影子定位的核心是一个正向模型和一个反向模型。正向模型给定日期、时间、经纬度和杆高，计算太阳在当地地平坐标系中的单位方向向量，再把杆顶沿反太阳方向投影到水平地面，得到影尖坐标。反向模型则把经纬度、杆高、坐标系旋转角以及可能的日期看作待估参数，使理论影尖轨迹尽可能接近附件中的观测轨迹。

本文的总体流程为：首先从题面和 Excel 附件中提取时间与影尖坐标；其次建立太阳位置和影子投影公式；然后对问题一直接正向计算影长，对问题二在已知日期下做非线性最小二乘反演，对问题三在全年日期网格上外层搜索日期、内层反演地点和杆高；最后对问题四给出视频抽帧和影尖检测流程，并说明缺失视频导致无法输出真实地点。所有关键结果均由 Python 脚本生成，并冻结到 \verb|results/frozen_numbers.json| 中。

\subsection{问题一分析}
问题一是典型的机理建模问题，已知日期、地点、杆高和时间区间，只需要根据太阳高度角计算影长。该问的难点不在优化，而在于正确处理北京时间、当地经度和太阳时之间的差异。若忽略经度修正，最短影长会被误认为发生在北京时间 12:00；实际上天安门经度为 $116.3914^\circ$E，相对北京时间标准经线 $120^\circ$E 有约 14.4 分钟差异，再加上时差方程修正，最短影长时间会偏离 12:00。

本文采用太阳赤纬、时角、太阳高度角模型。Baseline 为“正午高度角近似”，只判断影长量级；主模型进一步使用时差方程和经度修正，逐分钟计算太阳高度角并绘制曲线。

\subsection{问题二分析}
附件1已知日期为 2015 年 4 月 18 日，但题面未给出直杆高度，也未说明附件中的 $x,y$ 坐标轴是否对应正东、正北方向。因此，不能只用影子长度估算纬度，还必须利用二维影尖轨迹，同时允许坐标系相对于地理东--北坐标有一个未知旋转角。本文令未知参数为纬度 $\varphi$、经度 $\lambda$、杆高 $H$ 和旋转角 $\theta$，最小化所有时刻二维坐标残差平方和。

Baseline 采用影长最小时间估计经度、影长量级粗估纬度的方法。这一方法只能给出较粗位置，且无法处理坐标轴旋转。主模型使用完整二维轨迹，能同时解释影长变化和影子方向变化，因此更适合本问。

\subsection{问题三分析}
附件2和附件3缺少日期，未知量比问题二更多。若把日期当连续变量直接优化，容易陷入局部最优，也难以解释不同日期下的等价解。因此本文采用外层日期网格搜索：对 2015 年的每个候选日期，内层用问题二同样的最小二乘模型反演经纬度、杆高和旋转角；最后按 RMSE 排序并选出相隔较远的若干日期候选。

本问的重要特征是多解性。太阳赤纬在一年内关于夏至或冬至存在近似对称，某些春季与夏季、冬季与次年初日期可能产生非常接近的太阳高度角变化曲线。若观测时间段较短，仅凭一小时左右的影尖轨迹通常不能唯一确定日期。因此本文不只给一个最优解，而给出若干可能地点与日期，并用残差曲线解释多解来源。

\subsection{问题四分析}
问题四从结构化坐标数据扩展到视频数据。视频本身需要先经过图像处理，才能得到与附件1--3同类的影尖坐标序列。由于当前目录只有“附件4下载说明.doc”，没有视频文件，本文无法真实抽帧和测量影尖坐标。为保证学术诚信，本文不编造视频结果，而建立完整处理模型：抽取视频帧，识别杆底和影尖，利用已知 2 m 杆高完成像素到实际长度标定，再把得到的影尖坐标序列输入问题二或问题三模型。若日期已知，则只反演地点；若日期未知，则执行全年日期搜索并给出候选组合。

\section{模型假设与数据预处理}
\subsection{模型假设}
\begin{enumerate}[label=\textbf{假设\arabic*：}]
  \item 直杆竖直、地面水平。合理性：题面明确“直杆垂直于地面”，坐标系以杆底为原点且水平地面为 $xy$ 平面；作用：使影尖可由杆顶沿太阳光线投影到 $z=0$ 平面计算。
  \item 测量时间段内杆高、杆底位置和局部坐标系固定。合理性：附件数据来自同一固定直杆；作用：同一附件只需拟合一个杆高和一个平面旋转角。
  \item 附件坐标轴相对于地理东--北坐标可存在未知旋转。合理性：题面只定义 $xy$ 平面而未定义 $x$ 轴朝向；作用：模型引入旋转角 $\theta$，避免把坐标轴方向误当成真实地理方位。
  \item 同一附件内各影尖坐标测量误差独立且量级相近。合理性：附件数据保留四位小数，测量精度近似一致；作用：最小二乘目标函数使用未加权欧氏残差。
  \item 附件4视频缺失时不构造虚假影尖坐标。合理性：没有视频不能进行真实图像测量；作用：问题四只给出可执行模型流程和数据需求，不输出伪造地点。
\end{enumerate}

\subsection{数据来源与审计}
本文使用的原始数据位于题目目录中。题面文件为 \verb|CUMCM-2015-problem A-Chinese.doc|，附件1--3位于 \verb|附件1-3.xls|。附件4目录中仅有下载说明文件，给出网盘链接和赛题站链接，但没有视频本体。

\begin{table}[H]
\centering
\caption{附件数据审计表}
\begin{tabular}{lllll}
\toprule
附件 & 日期信息 & 时间范围 & 样本数 & 坐标字段\\
\midrule
附件1 & 2015-04-18 & 14:42--15:45 & 22 & $x,y$ 影尖坐标/m\\
附件2 & 未知 & 12:41--13:44 & 22 & $x,y$ 影尖坐标/m\\
附件3 & 未知 & 13:09--14:12 & 22 & $x,y$ 影尖坐标/m\\
附件4 & 未知/视频 & 视频文件缺失 & -- & 需由视频抽取得到\\
\bottomrule
\end{tabular}
\end{table}

数据预处理包括：读取 Excel 各工作表第 4 行起的时间与坐标；把北京时间字符串转换为 Python 时间对象；保留坐标原始单位“米”；检查空行和缺失值。附件1--3均为 22 条有效记录，无坐标缺失。由于数据点少且题目要求定位，未删除任何观测点，而是在模型检验中报告逐点残差。

\section{符号说明}
\begin{longtable}{p{0.18\textwidth}p{0.62\textwidth}p{0.12\textwidth}}
\caption{主要符号说明}\\
\toprule
符号 & 含义 & 单位\\
\midrule
\endfirsthead
\toprule
符号 & 含义 & 单位\\
\midrule
\endhead
$\varphi$ & 观测地点纬度，北纬为正 & $^\circ$\\
$\lambda$ & 观测地点经度，东经为正 & $^\circ$\\
$H$ & 直杆高度 & m\\
$n$ & 日期在一年中的日序 & d\\
$\delta$ & 太阳赤纬 & rad\\
$E_t$ & 时差方程修正量 & min\\
$T$ & 北京时间对应的钟表分钟数 & min\\
$\omega$ & 太阳时角 & rad\\
$\alpha$ & 太阳高度角 & rad\\
$\mathbf{s}$ & 太阳单位方向向量在东、北、天顶方向的分量 & --\\
$\mathbf{p}_i$ & 第 $i$ 个时刻理论影尖坐标 & m\\
$\mathbf{q}_i$ & 第 $i$ 个时刻附件观测影尖坐标 & m\\
$\theta$ & 附件坐标系相对东--北坐标的旋转角 & rad\\
$J$ & 最小二乘目标函数 & m$^2$\\
RMSE & 坐标拟合均方根误差 & m\\
\bottomrule
\end{longtable}

\section{模型建立与求解}
\subsection{太阳位置与影子投影的正向模型}
设观测点纬度为 $\varphi$，日期日序为 $n$。太阳赤纬采用 Cooper 近似公式
\begin{equation}
\delta=23.45^\circ\sin\left(\frac{360^\circ(284+n)}{365}\right).
\end{equation}
考虑时差方程
\begin{equation}
B=\frac{360^\circ(n-81)}{364},\quad E_t=9.87\sin 2B-7.53\cos B-1.5\sin B.
\end{equation}
北京时间以 $120^\circ$E 为标准经线，若钟表时间折算为分钟数 $T$，则真太阳时分钟数可写为
\begin{equation}
T_s=T+E_t+4(\lambda-120).
\end{equation}
对应时角为
\begin{equation}
\omega=15^\circ\left(\frac{T_s}{60}-12\right)=\frac{T_s}{4}-180^\circ.
\end{equation}
太阳高度角满足
\begin{equation}
\sin\alpha=\sin\varphi\sin\delta+\cos\varphi\cos\delta\cos\omega.
\end{equation}
在局部东--北--天顶坐标系中，太阳方向单位向量可写为
\begin{equation}
\mathbf{s}=\left(-\cos\delta\sin\omega,\ \cos\varphi\sin\delta-\sin\varphi\cos\delta\cos\omega,\ \sin\alpha\right)^T .
\end{equation}
若直杆高度为 $H$，杆顶为 $(0,0,H)$，沿反太阳方向投影到地面 $z=0$，得到地理东--北坐标下的影尖位置
\begin{equation}
\mathbf{r}(t)=-H\frac{(s_E(t),s_N(t))^T}{s_U(t)}.
\end{equation}
当只关心影长时，
\begin{equation}
L(t)=\|\mathbf{r}(t)\|=\frac{H}{\tan\alpha(t)}.
\end{equation}
该模型直接回答问题一，也作为问题二、三反演模型的基础。

\subsection{问题一：天安门影长曲线}
将天安门广场坐标 $\varphi=39^\circ54'26''$、$\lambda=116^\circ23'29''$，日期 2015 年 10 月 22 日，杆高 $H=3$ m 代入式(1)--(8)，逐分钟计算 9:00--15:00 的影长。计算得到 9:00 影长 6.8841 m，15:00 影长 7.0365 m，最短影长出现在 11:58，长度 3.8411 m。

\begin{figure}[H]
\centering
\includegraphics[width=0.86\textwidth]{\figpath 问题1_天安门影长曲线.png}
\caption{2015年10月22日天安门广场3m直杆影长变化曲线}
\end{figure}

由图1可见，上午 9:00 后太阳高度角逐渐增大，影长随之缩短；接近当地真太阳正午时达到最短；之后太阳高度角减小，影长再次增长。曲线最低点并非北京时间 12:00，而是约 11:58，这是经度修正和时差方程共同作用的结果。9:00 和 15:00 的影长接近但不完全对称，也反映了太阳赤纬和时差修正使钟表时间相对于太阳时存在偏移。

\subsection{问题二：已知日期下的地点反演模型}
附件坐标系可能与地理东--北坐标不一致，设旋转矩阵为
\begin{equation}
R(\theta)=\begin{pmatrix}\cos\theta&-\sin\theta\\\sin\theta&\cos\theta\end{pmatrix}.
\end{equation}
给定日期 $d$ 和时刻序列 $t_i$，理论影尖坐标为
\begin{equation}
\mathbf{p}_i(\varphi,\lambda,H,\theta)=R(\theta)\mathbf{r}(t_i;\varphi,\lambda,H,d).
\end{equation}
观测坐标为 $\mathbf{q}_i=(x_i,y_i)^T$，建立目标函数
\begin{equation}
\min_{\varphi,\lambda,H,\theta}J=\sum_{i=1}^{N}\left\|\mathbf{p}_i-\mathbf{q}_i\right\|_2^2,
\quad -65^\circ\le\varphi\le65^\circ,
\quad 70^\circ\le\lambda\le150^\circ,
\quad H>0.
\end{equation}
为保证 $H>0$，程序中实际优化变量为 $\log H$。求解采用多初值非线性最小二乘：先在纬度、经度、旋转角上设置粗网格初值，再用 \verb|scipy.optimize.least_squares| 局部收敛，最终取 RMSE 最小的解。

附件1日期为 2015 年 4 月 18 日。求解结果见表3。结果显示主要候选地点为北纬 18.3361$^\circ$、东经 109.7018$^\circ$，拟合杆高为 1.9640 m，旋转角为 14.9293$^\circ$。该经纬度位于海南岛附近，从太阳高度角和下午影子变化方向看具有现实可解释性。

\begin{table}[H]
\centering
\caption{附件1地点反演结果}
\begin{tabular}{llllll}
\toprule
日期 & 纬度/$^\circ$ & 经度/$^\circ$ & 杆高/m & 旋转角/$^\circ$ & RMSE/m\\
\midrule
2015-04-18 & 18.3361 & 109.7018 & 1.9640 & 14.9293 & 0.000251\\
\bottomrule
\end{tabular}
\end{table}

\begin{figure}[H]
\centering
\includegraphics[width=0.68\textwidth]{\figpath 问题2_附件1轨迹拟合.png}
\caption{附件1观测影尖轨迹与理论轨迹拟合}
\end{figure}

图2中观测点和理论点几乎重合，说明模型不仅拟合了影长，而且拟合了二维方向变化。与只用影长最小值的 baseline 相比，二维轨迹模型充分利用了 22 个观测点，能够估计坐标旋转角，并减少经纬度误判。

\begin{figure}[H]
\centering
\includegraphics[width=0.82\textwidth]{\figpath 问题2_附件1残差.png}
\caption{附件1逐时刻拟合残差}
\end{figure}

图3显示逐点残差均在毫米量级，最大残差为 0.000790 m。残差没有出现随时间单调增大的系统偏差，说明太阳位置公式、时间修正和坐标旋转的组合能够解释附件1的主要变化。Bootstrap 扰动实验中，纬度标准差约 0.0359$^\circ$、经度标准差约 0.0945$^\circ$，因此在题目数据精度下，该地点候选具有较强稳定性。

\subsection{问题三：未知日期下的地点--日期联合反演模型}
当日期未知时，令 $d$ 在 2015 年全年取值。对每个候选日期 $d$，求解
\begin{equation}
F(d)=\min_{\varphi,\lambda,H,\theta}\sum_{i=1}^{N}\left\|R(\theta)\mathbf{r}(t_i;\varphi,\lambda,H,d)-\mathbf{q}_i\right\|_2^2.
\end{equation}
日期--地点候选按
\begin{equation}
\mathrm{RMSE}(d)=\sqrt{F(d)/(2N)}
\end{equation}
排序。由于近邻日期通常给出非常接近的解，本文在输出候选时要求候选日期相隔至少约 12 天，以避免同一局部谷底被重复列出。

\subsubsection{附件2结果}
附件2最优候选为 2015 年 5 月 23 日、北纬 39.7296$^\circ$、东经 78.8739$^\circ$，拟合杆高 1.9936 m，RMSE 为 $3.4\times10^{-5}$ m。另一个几乎等价的候选为 2015 年 7 月 20 日、北纬 39.8040$^\circ$、东经 81.2112$^\circ$，RMSE 同为 $3.4\times10^{-5}$ m。两个日期分别位于夏至前后，太阳赤纬接近，因而在短时间观测窗口中产生几乎相同的影尖轨迹。

\begin{table}[H]
\centering
\caption{附件2若干可能日期与地点}
\begin{tabular}{llllll}
\toprule
序号 & 日期 & 纬度/$^\circ$ & 经度/$^\circ$ & 杆高/m & RMSE/m\\
\midrule
1 & 2015-05-23 & 39.7296 & 78.8739 & 1.9936 & 0.000034\\
2 & 2015-07-20 & 39.8040 & 81.2112 & 1.9954 & 0.000034\\
3 & 2015-06-04 & 41.1706 & 79.0361 & 2.0308 & 0.000162\\
4 & 2015-07-08 & 41.2148 & 80.6733 & 2.0320 & 0.000167\\
5 & 2015-08-01 & 37.6860 & 81.5377 & 1.9480 & 0.000231\\
6 & 2015-05-11 & 37.5830 & 79.1261 & 1.9460 & 0.000242\\
\bottomrule
\end{tabular}
\end{table}

\begin{figure}[H]
\centering
\includegraphics[width=0.86\textwidth]{\figpath 问题3_附件2_日期搜索残差.png}
\caption{附件2全年日期搜索残差曲线}
\end{figure}

图4中可以看到残差在 5 月下旬和 7 月下旬附近形成两个主要谷底。这说明在只有 63 分钟观测、且坐标轴旋转和杆高均未知的条件下，日期并不唯一。本文因此给出“若干可能地点与日期”，而不是强行给出唯一解。

\begin{figure}[H]
\centering
\includegraphics[width=0.68\textwidth]{\figpath 问题3_附件2_最佳轨迹拟合.png}
\caption{附件2最佳候选轨迹拟合}
\end{figure}

图5显示附件2最佳候选下理论轨迹与观测轨迹高度一致。由于 RMSE 仅为 $3.4\times10^{-5}$ m，坐标残差远小于附件坐标保留精度，表明模型在该候选下几乎完全重构了影尖轨迹。

\subsubsection{附件3结果}
附件3最优候选为 2015 年 11 月 16 日、北纬 30.4121$^\circ$、东经 106.5854$^\circ$，拟合杆高 2.9329 m，RMSE 为 0.000204 m。另一个近似等价候选为 2015 年 1 月 24 日、北纬 30.2418$^\circ$、东经 113.2011$^\circ$，RMSE 为 0.000205 m。二者分别位于冬至前后，太阳赤纬相近，同样体现出未知日期下的多解性。

\begin{table}[H]
\centering
\caption{附件3若干可能日期与地点}
\begin{tabular}{llllll}
\toprule
序号 & 日期 & 纬度/$^\circ$ & 经度/$^\circ$ & 杆高/m & RMSE/m\\
\midrule
1 & 2015-11-16 & 30.4121 & 106.5854 & 2.9329 & 0.000204\\
2 & 2015-01-24 & 30.2418 & 113.2011 & 2.9391 & 0.000205\\
3 & 2015-11-28 & 26.9038 & 107.0790 & 3.0660 & 0.000463\\
4 & 2015-01-12 & 26.7858 & 111.8606 & 3.0707 & 0.000475\\
5 & 2015-02-05 & 34.8073 & 114.3517 & 2.7816 & 0.000609\\
6 & 2015-11-04 & 35.0164 & 106.7637 & 2.7748 & 0.000636\\
\bottomrule
\end{tabular}
\end{table}

\begin{figure}[H]
\centering
\includegraphics[width=0.86\textwidth]{\figpath 问题3_附件3_日期搜索残差.png}
\caption{附件3全年日期搜索残差曲线}
\end{figure}

图6表明附件3残差谷底位于年初和年末两段，这与冬至前后太阳赤纬对称性一致。附件3影长明显较长，拟合杆高约 2.93 m，说明其直杆高度可能不同于附件2，而题目并未限定附件之间杆高相同。

\begin{figure}[H]
\centering
\includegraphics[width=0.68\textwidth]{\figpath 问题3_附件3_最佳轨迹拟合.png}
\caption{附件3最佳候选轨迹拟合}
\end{figure}

图7中附件3观测轨迹和拟合轨迹同样基本重合。相较附件2，附件3最优 RMSE 略大，但仍在 0.3 mm 以内，远低于一般视频或人工测量误差。因此，附件3结果的主要不确定性不是坐标拟合误差，而是日期多解性和地理先验不足。

\subsection{问题四：视频影子定位模型}
若获得附件4视频，本文建议按以下步骤处理。
\begin{enumerate}
  \item \textbf{视频抽帧：}按固定时间间隔抽取帧，记录每帧对应北京时间。若视频无绝对时间戳，则以首帧时间为未知量，可在外层一并搜索或由题目元数据确定。
  \item \textbf{杆底与影尖检测：}对每帧进行灰度化、背景/阴影分割和边缘检测；杆底可通过直杆与地面的交点固定，影尖取阴影区域最远端或经人工少量校正。
  \item \textbf{像素--米标定：}题目给出杆高为 2 m，可利用图像中直杆像素高度或相机投影近似建立局部比例尺。若相机俯仰显著，应先做地面单应性校正。
  \item \textbf{轨迹反演：}把每帧影尖相对杆底的地面坐标输入式(10)--(12)。若日期已知，按问题二反演地点；若日期未知，按问题三搜索日期--地点候选。
  \item \textbf{结果检验：}用理论轨迹重投影到视频帧，计算像素残差；检查影尖轨迹是否随时间平滑，排除遮挡和检测错误帧。
\end{enumerate}

当前目录缺少视频文件，且自动访问下载站点失败。本文在支撑材料 \verb|qa/preflight_report.md| 中记录了该事实。因此，问题四不输出伪造拍摄地点；若用户补充视频，可直接运行上述流程，且前 3 问代码已提供地点/日期反演核心函数。

\section{模型检验、对比与敏感性分析}
\subsection{Baseline 对比}
问题一的 baseline 是忽略时差方程和经度修正的正午近似，只能给出“中午最短”的定性判断；主模型给出了 11:58 的具体最短时刻和逐分钟曲线。问题二、三的 baseline 是利用影长最短时刻和影长量级粗定位，但无法处理未知坐标轴方向，也不能解释二维轨迹。主模型通过引入旋转角，把完整轨迹用于拟合，残差达到毫米甚至更小，显著优于 baseline 的粗定位能力。

\begin{table}[H]
\centering
\caption{主模型与baseline能力对比}
\begin{tabular}{p{0.18\textwidth}p{0.35\textwidth}p{0.35\textwidth}}
\toprule
问题 & Baseline & 主模型优势\\
\midrule
问题一 & 正午高度角近似 & 加入经度和时差方程，得到分钟级曲线与极值\\
问题二 & 影长最小值粗定位 & 拟合二维轨迹、杆高和旋转角，RMSE为0.000251 m\\
问题三 & 逐日影长相似度比较 & 日期网格+地点联合反演，输出多组可解释候选\\
问题四 & 少量人工标注影长 & 完整视频抽帧、影尖检测、重投影残差检验\\
\bottomrule
\end{tabular}
\end{table}

\subsection{残差与约束检验}
附件1最大残差 0.000790 m，附件2最佳候选 RMSE 为 $3.4\times10^{-5}$ m，附件3最佳候选 RMSE 为 0.000204 m。所有最优解均满足经纬度边界、杆高为正和太阳高度角为正等约束。残差图未显示明显单调偏差，说明没有出现由时间读取错误、经纬度修正方向错误或坐标轴旋转遗漏导致的系统误差。

\subsection{敏感性分析}
对附件1，本文在日期固定值附近做 $\pm2$ 天扰动并重新拟合地点。若强行使用错误日期，模型仍可通过调整经纬度和杆高降低残差，但 RMSE 会增大，且地点发生可观察偏移。这说明问题二中“日期已知”是重要信息；如果日期不确定，应采用问题三的日期网格搜索，而不能把某一日期的反演结果视为唯一真实地点。

对问题三，日期搜索残差曲线本身就是关于日期的敏感性分析。附件2在 5 月 23 日和 7 月 20 日附近存在几乎等深谷底，附件3在 11 月 16 日和 1 月 24 日附近存在近似等深谷底。这说明短时影子轨迹对“太阳赤纬”敏感，但对赤纬相近的两个日期不敏感。因此，若要消除多解，需要更长时间跨度、已知杆高、坐标轴朝向、地理范围先验或视频元数据等额外信息。

\section{模型评价、改进与推广}
\subsection{模型优点}
第一，模型具有明确物理含义。每一步都来自太阳位置和几何投影，而不是黑箱拟合，因此结果可以通过太阳高度角、时角和影子方向解释。第二，模型能处理题目中的实际不确定性。附件只给局部 $x,y$ 坐标而没有说明坐标轴方向，本文引入旋转角，使模型不依赖额外坐标轴假设。第三，结果可复现。全部代码、图表、结果表和冻结数字均保存在支撑材料中，论文中的关键数字均来自 \verb|frozen_numbers.json|。

\subsection{模型局限}
第一，太阳位置公式采用竞赛建模常用近似，未使用高精度天文历表；在本题数据精度下残差已很小，但若用于严肃天文测量，应替换为更高精度的太阳历算法。第二，未知日期问题存在固有多解性。短时间影尖轨迹无法唯一地区分太阳赤纬相近的两个日期，因此本文只能给出若干候选。第三，问题四缺少视频文件，本文无法完成真实图像测量，只能给出可执行流程和模型框架。

\subsection{改进方向}
若需要提高定位唯一性，可以引入更多先验：例如已知杆高、拍摄大致城市范围、坐标轴朝向、视频拍摄时间戳或更长的观测时段。对于视频问题，可使用交互式标注少量关键帧与自动影尖检测相结合的方式，提高影尖坐标质量；也可以使用相机标定和地面单应性变换减少透视误差。对于太阳位置计算，可用 NOAA SPA 或天文历表替换本文近似公式。

\subsection{推广应用}
本文方法不仅适用于直杆影子定位，也可推广到建筑物、路灯、旗杆等竖直物体的影子定位。在已知物体高度和视频时间的情况下，可用于地理取证；在已知地理位置的情况下，也可以反向校验视频拍摄时间。对于无人机影像和街景图像，只要能提取物体底点和影尖点，本文的投影和反演框架仍然适用。

\section{结论}
本文完成了2015A太阳影子定位题的建模、求解、检验和论文支撑材料整理。主要结论如下：
\begin{enumerate}
  \item 天安门广场 2015 年 10 月 22 日三米直杆影长在 9:00--15:00 内先减小后增大；9:00 为 6.8841 m，15:00 为 7.0365 m，11:58 达到最小值 3.8411 m。
  \item 附件1在已知日期 2015 年 4 月 18 日下，主要可能地点为北纬 18.3361$^\circ$、东经 109.7018$^\circ$，拟合杆高 1.9640 m，RMSE 为 0.000251 m。
  \item 附件2未知日期下，优先候选为 2015 年 5 月 23 日、北纬 39.7296$^\circ$、东经 78.8739$^\circ$；2015 年 7 月 20 日、北纬 39.8040$^\circ$、东经 81.2112$^\circ$为近似等价候选。
  \item 附件3未知日期下，优先候选为 2015 年 11 月 16 日、北纬 30.4121$^\circ$、东经 106.5854$^\circ$；2015 年 1 月 24 日、北纬 30.2418$^\circ$、东经 113.2011$^\circ$为近似等价候选。
  \item 附件4视频文件当前缺失，本文不伪造地点；已给出抽帧、影尖检测、像素标定和地点/日期反演的完整模型流程。
\end{enumerate}

\section*{参考文献}
\addcontentsline{toc}{section}{参考文献}
\begin{enumerate}[label={[\arabic*]}]
\item Duffie J A, Beckman W A. Solar Engineering of Thermal Processes[M]. Hoboken: Wiley, 2013.
\item Reda I, Andreas A. Solar position algorithm for solar radiation applications[J]. Solar Energy, 2004, 76(5): 577--589.
\item Meeus J. Astronomical Algorithms[M]. Richmond: Willmann-Bell, 1998.
\item Spencer J W. Fourier series representation of the position of the sun[J]. Search, 1971, 2(5): 172.
\item 全国大学生数学建模竞赛组委会. 2015高教社杯全国大学生数学建模竞赛A题：太阳影子定位[Z]. 2015.
\item Virtanen P, Gommers R, Oliphant T E, et al. SciPy 1.0: fundamental algorithms for scientific computing in Python[J]. Nature Methods, 2020, 17: 261--272.
\end{enumerate}

\appendix
\section{支撑材料说明}
支撑材料按题目目录原地建立，主要包括：\verb|代码/2015A_shadow_modeling.py| 为完整求解脚本；\verb|图表/| 保存所有论文图；\verb|tables/| 保存逐分钟影长、拟合轨迹、候选结果和敏感性表；\verb|results/frozen_numbers.json| 保存论文引用的冻结数字；\verb|contracts/| 保存题意、模型路线和结果契约；\verb|qa/| 保存预检和质量门控记录。

\section{代码入口与运行环境}
代码使用 Python 3.11/3.12，依赖 \verb|numpy|、\verb|pandas|、\verb|scipy|、\verb|matplotlib|、\verb|xlrd|。在支撑材料目录下可运行：
\begin{verbatim}
python 代码/2015A_shadow_modeling.py
\end{verbatim}
脚本会自动读取 \verb|数据/附件1-3.xls|，并重新生成全部图表、结果表和冻结数字。

\section{三层质量门控摘要}
\begin{itemize}
  \item L1 建模合理性：通过。每问均有题目输出映射，假设与变量用于模型，baseline 与主模型差异明确。
  \item L2 求解正确性：通过。数据审计完成，代码可复现，残差和候选结果已冻结，图表由代码生成。
  \item L3 论文质量：通过。摘要含关键数字，图表均有编号和解释，LaTeX 编译并核验页数，支撑材料已打包。
\end{itemize}

\end{document}
